\label{sec:results}

\subsection{Diagnostic: map-free generation fails at corridor-level realism}
\label{sec:results-diagnostic}

We first establish a diagnostic finding under the redefined notion of realism (Sec.~\ref{sec:intro}): without an explicit structural mechanism, map-free continuous generation struggles to produce trip-scale routes that align with road corridors.
Figure~\ref{fig:diagnosis-mapfree} visualizes a representative OD pair with multiple GT trajectories that form two distinct corridors. In contrast, an L2 regressor collapses to a straight-line-like mode; end-to-end diffusion tends to blur around the origin; and map-free CascadeTraj produces zigzag trajectories that frequently cross non-road areas.
Moreover, even when GT corridors are approximately balanced, the generated samples can be heavily skewed toward one corridor, indicating a biased decision distribution.

For the case shown in Fig.~\ref{fig:diagnosis-mapfree}, GT corridors are close to balanced ($10$ vs.\ $12$ trajectories), while the map-free CascadeTraj samples are skewed ($16$ vs.\ $4$ out of $20$), suggesting that decision-stage uncertainty is not well-calibrated.

\begin{figure}[t]
  \centering
  \maybeincludegraphics[width=\linewidth]{figures/fig1_diagnosis_mapfree.png}
  \caption{\textbf{Map-free generation fails at corridor-level realism.}
  Grey shading indicates an OSM-derived road mask (road probability $>0.5$) used only for visualization; the circle/square mark origin/destination. GT trajectories exhibit multi-modal corridor structure, while map-free baselines either collapse (L2), diverge/blur (end-to-end diffusion), or produce off-road zigzags (map-free CascadeTraj). All methods use the same OD pair and $K{=}20$ samples.}
  \label{fig:diagnosis-mapfree}
\end{figure}

\begin{figure}[t]
  \centering
  \maybeincludegraphics[width=0.9\linewidth]{figures/fig2_gt_corridors.png}
  \caption{\textbf{Ground-truth corridor multi-modality.}
  For the same OD, GT trajectories cluster into two distinct corridors (colors), motivating corridor-level modeling and distributional evaluation.}
  \label{fig:gt-corridors}
\end{figure}

\subsection{Waypoint audit: soft priors help but do not solve the measure-zero problem}
\label{sec:results-waypoint-audit}

To quantify the decision-stage failure mode, we audit decision-stage waypoints against GT corridors.
We measure the distance from generated waypoints to the GT polyline in \emph{grid units} (lower is better).
Table~\ref{tab:wp-audit} compares two map-free decision settings: OD-only conditioning and tier-road conditioning.
Tier-road reduces the error substantially, but the median distance remains far beyond the scale of road width under our grid, indicating that soft map priors alone do not yield corridor-level convergence.

\begin{table}[t]
  \centering
  \caption{\textbf{Decision waypoint distance to GT corridor (grid units).}
  Even with tier-road conditioning, waypoints remain far from GT corridors.}
  \label{tab:wp-audit}
  \begin{tabular}{lccc}
    \toprule
    Setting & Mean & P50 & P90 \\
    \midrule
    OD-only (map-free) & 59.37 & 49.51 & 125.19 \\
    + tier-road (soft prior) & 35.00 & 24.70 & 83.03 \\
    \bottomrule
  \end{tabular}
\end{table}

\begin{figure}[t]
  \centering
  \begin{subfigure}[t]{0.49\linewidth}
    \centering
    \maybeincludegraphics[width=\linewidth]{figures/fig_waypoint_audit_odonly.png}
    \caption{OD-only decision vs.\ GT}
  \end{subfigure}
  \hfill
  \begin{subfigure}[t]{0.49\linewidth}
    \centering
    \maybeincludegraphics[width=\linewidth]{figures/fig_waypoint_scatter_tierroad.png}
    \caption{Tier-road conditioned decision}
  \end{subfigure}
  \caption{\textbf{Decision-stage waypoints do not converge to roads.}
  Grey shading indicates the road mask; blue points are sampled decision waypoints, and black dots are waypoints extracted from GT trajectories along the GT corridor polyline. Tier-road conditioning reduces dispersion but waypoints remain far from corridors (Table~\ref{tab:wp-audit}).}
  \label{fig:waypoint-audit}
\end{figure}

\subsection{Implication: corridor choice requires discrete structure and semantic conditioning}
\label{sec:results-roadaware}

The diagnostic evidence supports two design principles for behaviorally realistic route generation:
\textbf{(i) Structural mechanism}: corridor choice should be a discrete decision on the road graph, not a continuous optimization toward a measure-zero manifold.
\textbf{(ii) Behavioral mechanism}: corridor probabilities must change with temporal--spatial context---not randomly, but in response to semantics (POI activity, road hierarchy) and time (rush hour vs.\ midnight).
CascadeTraj instantiates both: the road graph provides structural validity, while semantic/temporal conditioning provides behavioral grounding. This turns diversity from indiscriminate randomness into context-dependent behavioral variation.
% =====================================================================
% LEGACY (2026-01-14)
% This file is NOT included by `essay_icml_cascadetraj/main.tex`.
% It is kept only for historical reference. Do not edit for the ICML
% routegen mainline. See `sections/06_results.tex` + `sections/07_discussion.tex`.
% =====================================================================
